\documentclass[11pt]{article}

\usepackage[T1]{fontenc}
\usepackage{lmodern}
\usepackage{amsmath,amssymb,amsthm,mathtools}
\usepackage{booktabs,array}
\usepackage{enumitem}
\usepackage[a4paper,margin=27mm]{geometry}
\usepackage[colorlinks=true,linkcolor=blue,citecolor=blue,urlcolor=blue]{hyperref}
\setlength{\emergencystretch}{2em}
\numberwithin{equation}{section}

\newtheorem{theorem}{Theorem}[section]
\newtheorem{proposition}[theorem]{Proposition}
\newtheorem{lemma}[theorem]{Lemma}
\newtheorem{corollary}[theorem]{Corollary}
\theoremstyle{definition}
\newtheorem{example}[theorem]{Example}
\theoremstyle{remark}
\newtheorem{remark}[theorem]{Remark}

\newcommand{\CC}{\mathbf C}
\newcommand{\ZZ}{\mathbf Z}
\newcommand{\PP}{\mathbf P}
\newcommand{\OO}{\mathcal O}
\newcommand{\omegao}{\omega_0}
\newcommand{\bfzero}{b_{f,0}}
\DeclareMathOperator{\ord}{ord}
\DeclareMathOperator{\Res}{Res}
\DeclareMathOperator{\divisor}{div}
\DeclareMathOperator{\Poles}{Poles}
\DeclareMathOperator{\mult}{mult}

\title{Uniform Bernstein--Sato Factors and Contact Roots\\
for Surface Germs with a Multiple Smooth Tangent Cone}
\author{YINGFENG JIANG}
\date{20 October 2026}

\begin{document}
\maketitle

\begin{abstract}
Let $f=g_e^r+h_D+R_{>D}$ be a convergent analytic germ in three
variables, where $D=er+k$, $e\ge3$, $r\ge2$, $k\ge1$, the plane
curve $C=V(g_e)$ is smooth, and $h_D|_C\not\equiv0$.
We prove that the local Bernstein--Sato polynomial contains the
uniform factor
$\prod_{j=0}^{D-1}(s+(3r+k)/(rD)+j/D)$, independently of the
contact data and the higher analytic terms. Holomorphic numerators
give a triangular set of exact roots containing this block.
For the two-layer
germs $g_e^r+h_D$, with $V(h_D)$ smooth at its intersections with
$C$, we determine a root for every nontransverse contact order
$q\ge2$ and prove that every pole of the local topological zeta
function is a root. An explicit infinite family has a contact root
outside the triangular set even though the corresponding topological pole
cancels. The proof combines compact residue periods with a cyclic
cover of $C$; ramification and Riemann--Roch supply the numerators
that fill the uniform factor.
\end{abstract}

\section{Introduction}\label{sec:introduction}

For surface germs whose first two homogeneous layers are
$g_e^r+h_D$, with a smooth projective curve $C=V(g_e)$,
we construct a uniform block of $D$ exact Bernstein--Sato
roots. This block depends only on $e,r,D$ and persists
under arbitrary higher analytic terms. Nontransverse
contacts give a second source of roots. In an explicit
infinite family, a contact root lies outside the entire
triangular root set even though its topological pole cancels.

Let $g_e,h_D\in\CC[x,y,z]$ be homogeneous of degrees $e,D$.
For integers $e,r,k$, set
\begin{equation}\label{eq:degrees}
 A=er,\qquad D=er+k,\qquad e\ge3,\quad r\ge2,\quad k\ge1.
\end{equation}
We assume that $g_e$ defines a smooth plane curve $C\subset\PP^2$
of degree $e$ and that $h_D|_C\not\equiv0$. The first theorem
applies to every convergent germ
\begin{equation}\label{eq:analytic-family}
 f=g_e^r+h_D+R_{>D},\qquad \ord_0R_{>D}>D.
\end{equation}
Write $\bfzero(s)$ for its monic local Bernstein--Sato polynomial:
it generates the polynomials $b(s)$ for which
$b(s)f^s=P(s)f^{s+1}$ with $P(s)\in\mathcal D_{\CC^3,0}[s]$.
Put
\begin{equation}\label{eq:alpha-R}
 \alpha_R=\frac{3r+k}{rD}.
\end{equation}

\begin{theorem}[Uniform factor and triangular root set]
\label{thm:uniform}
For every germ \eqref{eq:analytic-family},
\begin{equation}\label{eq:uniform-factor}
 \prod_{j=0}^{D-1}\left(s+\alpha_R+\frac jD\right)
 \ \mid\ \bfzero(s).
\end{equation}
More generally, for integers $\ell,j$,
\begin{equation}\label{eq:triangular-roots}
 \bfzero\!\left(-\alpha_R-\frac\ell r-\frac jD\right)=0
 \quad\text{if}\quad
 \begin{cases}
  0\le\ell\le r-2,\quad 0\le j<D,\\
  \ell/r+j/D<1.
 \end{cases}
\end{equation}
\end{theorem}

The factor in \eqref{eq:uniform-factor} depends only on $e,r,k$.
In particular, the theorem allows arbitrary zero multiplicities
of $h_D|_C$ and nonisolated germs. The triangular set
\eqref{eq:triangular-roots} contains the uniform block and
strictly enlarges it exactly when $r\ge3$ and $r\nmid k$.
It is obtained from the original
holomorphic numerators $g_e^\ell p_j$, where $p_j$ has degree $j$.
Different pairs $(\ell,j)$ may give the same root; throughout the
paper root sets are counted without repetition.

For the contact results, suppose that
\begin{equation}\label{eq:two-layer}
 f=g_e^r+h_D,\qquad H=V(h_D)\text{ is smooth at every point of }
 P=C\cap H.
\end{equation}
Smoothness here means that the local defining equation of
$H$ has nonzero differential at each point of $P$.
The set $P$ is finite. Let $q_p=I_p(C,H)$ be the actual
intersection multiplicity, and define
\begin{equation}\label{eq:contact-data}
 B=(r-1)D+k,\qquad \beta=3(r-1)+k,\qquad
 \alpha_q=\frac{q\beta+k}{qB}.
\end{equation}

\begin{theorem}[Contact roots]\label{thm:contact}
For every contact order $q=q_p\ge2$ in \eqref{eq:two-layer},
\begin{equation}\label{eq:contact-root}
 \bfzero(-\alpha_q)=0.
\end{equation}
\end{theorem}

The contact formula detects candidate exponents even after
they cancel from the topological zeta function. The following
family places one of these roots outside the triangular set.

\begin{theorem}[Contact roots after pole cancellation]
\label{thm:cancelled}
Suppose that the integers $r,e,k,q$ satisfy
\begin{equation}\label{eq:cancelled-family}
 r\ge3,\qquad k=r-1,\qquad e\ge3,\qquad
 D=r(e+1)-1=4q.
\end{equation}
There exist pairs $(C,H)$ as in \eqref{eq:two-layer} with one
contact of order $q$ and all other intersections transverse.
For every such pair,
\begin{equation}\label{eq:cancelled-conclusion}
 -\frac1q\notin\Poles Z_{\mathrm{top},0}(f;s),\qquad
 \bfzero(-1/q)=0.
\end{equation}
The root $-1/q$ lies outside the triangular set
\eqref{eq:triangular-roots}.
\end{theorem}

An infinite subfamily has $r=3$, $k=2$, $e=4m+2$ and
$q=3m+2$, for integers $m\ge1$. Its first value,
$(e,r,k,q)=(6,3,2,5)$, has at least $37$ distinct
Bernstein--Sato roots; see Example~\ref{ex:thirty-seven}.

The pole calculation gives the following application.

\begin{corollary}[Strong topological monodromy for the two-layer family]
\label{cor:strong}
Under \eqref{eq:two-layer}, the germ $f$ is reduced and has an
isolated critical point at the origin. Every pole of
$Z_{\mathrm{top},0}(f;s)$ is simple, and
\begin{equation}\label{eq:strong}
 s_0\in\Poles Z_{\mathrm{top},0}(f;s)
 \quad\Longrightarrow\quad \bfzero(s_0)=0.
\end{equation}
\end{corollary}

\subsection*{Period construction}

The uniform factor comes from a divisor lying over $C$. A cyclic
cover of $C$ separates its residue into an Euler form in the
fiber and a differential on a compact curve. The key step is to
choose $p_j$ so that the latter differential has a nonzero
period. At a zero of $h_D|_C$ of order $q_x$, its regularity
condition is
the explicit inequality
\[
 \ord_x(p_j|_C)\ge
 \left\lfloor\frac{(j+3-e)q_x}{D}\right\rfloor.
\]
The low degrees give holomorphic forms directly. For
$e\ge4$, Riemann--Roch fills the remaining degrees $j<D$.
For a cubic $C$, a degree-zero line bundle can require a
differential with simple poles and a nonzero residue period.
The corresponding original numerator is holomorphic.
These constructions fill the entire block in
\eqref{eq:uniform-factor}.

At a contact point, the prescribed residue is a two-dimensional
Euler form. In the exceptional case where its first parameter
is $a=1$, the contact root belongs to the curve's uniform block.
For all other parameters, explicit de Rham identities give a
nonzero period through a convergent positive integral or a
logarithmic residue.
Section~\ref{sec:periods} proves the period criterion,
Section~\ref{sec:uniform} constructs the uniform roots, and
Section~\ref{sec:contact} proves the contact theorem.
Sections~\ref{sec:poles} and \ref{sec:examples} give the
topological application and the cancellation examples.

\subsection*{Relation with earlier work}

Partial resolutions and their zeta functions are classical
tools for surface singularities; see Gusein-Zade, Luengo, and
Melle-Hern\'andez~\cite{GLM97}. Denef and Loeser's Newton
formula~\cite{DL92} supplies the local orbit calculation used
here, with the Jacobian of the first blow-up retained. Related
suspension and L\^e--Yomdin formulas are developed
in~\cite{AGGL26}. For the degree-gap-one family, the companion
paper~\cite{Jiang26} computes the topological and monodromy
invariants and reconstructs the contact orders from either.
The present paper constructs exact Bernstein--Sato roots for
arbitrary degree gap and a uniform factor in the broader
analytic family \eqref{eq:analytic-family}.

Our analytic input is the residue-period argument, in the form
of the all-numerator criterion of Esterov~\cite[Lemma~7.3]{Est26}.
Section~\ref{sec:periods} gives its divisor form with a
holomorphic numerator.

Blanco~\cite[Theorem~B and Proposition~C]{Bla26} constructs
geometric roots for plane curves on selected rupture divisors
from holomorphic numerators with $0\le v_i(g)<N_i$ and
residue conditions, and compares the resulting roots with
jumping numbers and motivic poles.
For our three-variable family, ramification of the cyclic curve
cover and Riemann--Roch supply original numerators and
nonzero periods for every $j<D$, producing the entire
uniform factor.

Examples of Bernstein--Sato roots absent from the motivic
zeta pole set occur in Bath and Veys~\cite[Example~2.13]{BV26}.
Theorem~\ref{thm:cancelled} supplies an explicit infinite
contact family, with a nonzero specified residue period and
a root outside the triangular set \eqref{eq:triangular-roots}.

\section{Compact residue periods and exact roots}\label{sec:periods}

A compact period of the prescribed normal residue detects
an exact Bernstein--Sato root. The criterion below keeps
track of the numerator's vanishing order; Euler forms then
supply the periods. All cycles have complex coefficients
and finite, hence compact, support.

\subsection{A divisor criterion with a holomorphic numerator}

Let $\pi:Y\to(\CC^n,0)$ be an embedded resolution, an
isomorphism outside $f^{-1}(0)$, and let $E$ be a prime divisor
with $\pi(E)=\{0\}$. On a smooth open subset $E^\circ$ meeting
no other component of the total transform, write
\begin{equation}\label{eq:divisor-model}
 f\circ\pi=t^NU(t,y),\qquad
 \pi^*\omegao=t^{\nu-1}J(t,y)\,dt\wedge dy,
\end{equation}
where $\omegao=dx_1\wedge\cdots\wedge dx_n$ and $U,J$ are
units. In the normal line bundle of $E$, define
\begin{equation}\label{eq:normal-cover}
 \widehat E^\circ=\{(y,\tau):\tau^NU(0,y)=1\}.
\end{equation}
This definition is intrinsic: if $t_b=c_{ba}t_a+O(t_a^2)$,
then $\tau_b=c_{ba}\tau_a$ and
$U_b(0,y)=c_{ba}^{-N}U_a(0,y)$.

\begin{proposition}[Numerator period criterion]\label{prop:criterion}
Let $p\in\OO_{\CC^n,0}$ and write
$p\circ\pi=t^mP(t,y)$, with $0\le m<N$. If the form
\begin{equation}\label{eq:specified-residue}
 \rho_p=\tau^{\nu+m}J(0,y)P(0,y)\,dy
\end{equation}
has a nonzero period on a compact $(n-1)$-cycle in
$\widehat E^\circ$, then
\begin{equation}\label{eq:criterion-root}
 b_{f,0}\!\left(-\frac{\nu+m}{N}\right)=0.
\end{equation}
\end{proposition}

\begin{proof}
Comparing the leading terms of $\pi^*(p\omegao)$ on
overlapping charts gives
\[
 J_b(0,y_b)P_b(0,y_b)\,dy_b
   =c_{ba}^{-(\nu+m)}J_a(0,y_a)P_a(0,y_a)\,dy_a.
\]
Together with $\tau_b=c_{ba}\tau_a$, this proves that
\eqref{eq:specified-residue} is a global form on
$\widehat E^\circ$.
Fix a cycle $\Gamma$ with $c=\int_\Gamma\rho_p\ne0$.
Its support is covered by finitely many charts
\eqref{eq:divisor-model}. Choose a smooth tubular map
$\Phi$ whose normal derivative along $E$ is complex linear,
and put $\epsilon=z^{1/N}>0$. Uniformly near this support,
\[
 \epsilon^{-N}(f\circ\pi)
       (\Phi(y,\epsilon s\tau))=s^N+O(\epsilon).
\]
At $(s,\epsilon)=(1,0)$, the real two-dimensional derivative
in the complex variable $s$ is invertible. The implicit function
theorem gives a unique smooth solution
$s_\epsilon=1+O(\epsilon)$ for which this expression equals
one. Uniqueness makes the local solutions agree. The resulting
map to $(f\circ\pi)^{-1}(z)$ satisfies
\[
 t=\epsilon\tau+O(\epsilon^2),\qquad
 y_{\epsilon}=y+O(\epsilon)
\]
in tangential $C^1$, uniformly on the fixed support. It sends
$\Gamma$ to a closed cycle; pushing it by $\pi$ gives a cycle
$\gamma_z\subset f^{-1}(z)$ whose support tends uniformly to
the origin.

The relative Gelfand--Leray form in these coordinates is
\[
 \frac{\pi^*\omegao}{d(f\circ\pi)}
 =\frac{t^{\nu-N}J(t,y)}{NU(t,y)+t\partial_tU(t,y)}\,dy.
\]
Its denominator is bounded away from zero on the chosen
neighborhood. Substitution gives, with
$\alpha_p=(\nu+m)/N$ and $\alpha_0=\nu/N$,
\begin{equation}\label{eq:period-asymptotics}
 \begin{split}
 I_p(z):=\int_{\gamma_z}\frac{p\omegao}{df}
   &=\frac1N z^{\alpha_p-1}
        \bigl(c+O(z^{1/N})\bigr),\\
 I_h(z):=\int_{\gamma_z}\frac{h\omegao}{df}
   &=O_h(z^{\alpha_0-1})
       \quad(h\in\OO_{\CC^n,0}).
 \end{split}
\end{equation}
The second estimate applies to every original holomorphic
numerator $h$.

Take a Bernstein identity $b(s)f^s=\mathcal P(s)f^{s+1}$. Multiplying
by $p\omegao$ and integrating the derivatives in $\mathcal P(s)$ by
parts gives
\begin{equation}\label{eq:green-identity}
 b(s)p f^s\omegao=f^{s+1}h_s\omegao+d\Xi_s,
 \qquad h_s=\sum_{a=0}^L s^a h_a,
\end{equation}
Here $h_s=\mathcal P(s)^*p$, where the star is the formal adjoint
with respect to $\omegao$; in particular, every $h_a$ is
holomorphic in the original space.
The form $\Xi_s$ is a finite sum of terms
$f^{s+1-b}\eta_b(s)$ with holomorphic $(n-1)$-forms
$\eta_b(s)$ polynomial in $s$.
For such a form, $df\wedge\eta_b=h\omegao$ with $h$
holomorphic, so its fiber integral has the second bound in
\eqref{eq:period-asymptotics}.

Integrate \eqref{eq:green-identity} over the chain swept out
by $\gamma_z$, $0<z\le z_0$. For $\Re s$ sufficiently
large the lower boundary tends to zero, and
\begin{equation}\label{eq:mellin-identity}
 b(s)M_p(s)=\sum_{a=0}^L s^aM_{h_a}(s+1)+B(s),
 \qquad M_h(s)=\int_0^{z_0}z^sI_h(z)\,dz,
\end{equation}
with $B(s)$ entire. Both sides are holomorphic for
$\Re s>-\alpha_p$, so the identity holds on that half-plane.
The right side is moreover holomorphic for
$\Re s>-\alpha_0-1$. Since $m<N$, it is holomorphic near
$s=-\alpha_p$. The first asymptotic gives
\[
 \lim_{u\downarrow0}uM_p(-\alpha_p+u)=c/N\ne0.
\]
Multiply \eqref{eq:mellin-identity} by $u$ and let
$s=-\alpha_p+u$ tend to $-\alpha_p$. This proves
$b(-\alpha_p)=0$, in particular for $b=\bfzero$.
\end{proof}

\subsection{Euler forms with compact periods}

All Euler parameters below are rational. A finite Kummer
cover makes the forms single-valued. The following elementary
regularization will also specify the cycles.

\begin{lemma}[Positive Euler periods]\label{lem:positive}
For $m=1,2$, let
\[
 M_m=\{z_1\cdots z_{m+1}\ne0,\quad
              z_1+\cdots+z_{m+1}=1\}.
\]
If every $\lambda_i$ is positive and nonintegral, then
$\prod_i z_i^{\lambda_i-1}dz_1\wedge\cdots\wedge dz_m$
has a nonzero compact period on its finite Kummer cover.
\end{lemma}

\begin{proof}
Truncate the real simplex $z_i>0$ near its faces and choose
the positive branch. A face has normal monodromy
$\zeta_i\ne1$. A small normal circle has twisted boundary
$\zeta_i-1$, so dividing by this scalar supplies a tube
that cancels the truncated boundary.
For $m=2$, the remaining corner boundaries are filled by
two-torus chains. More explicitly, at a corner put
$a=\zeta_i-1$, $b=\zeta_j-1$. The local boundary complex is
\[
 \CC\xrightarrow{c\mapsto(-bc,ac)}\CC^2
       \xrightarrow{(x,y)\mapsto ax+by}\CC.
\]
Every cycle $(x,y)$ in the middle is the image of $y/a$.
Thus the face tubes and corner chains close the truncated
simplex to a finite twisted cycle.

Changing the cutoff gives homologous cycles. Positivity of
the parameters makes the integrals over the added tubes
tend to zero with the cutoff. The cycle integral therefore
equals the strictly positive convergent simplex integral.
Finally, lift the chain to the finite cover and sum its lifts
with the character coefficients. This produces an ordinary
complex-coefficient cycle with a nonzero integral.
\end{proof}

Write $W=1-X-Y$ on $M_2$. The surface and curve Euler forms are
\begin{equation}\label{eq:euler-forms}
 \Phi_{a,b,c}=X^{a-1}Y^{b-1}W^{c-1}dX\wedge dY,
 \qquad
 \varphi_{a,b}=z^{a-1}(1-z)^{b-1}dz.
\end{equation}

\begin{lemma}[Euler nonvanishing]\label{lem:euler}
Each of the following forms has a nonzero compact period
on its finite Kummer cover:
\begin{enumerate}[label=\textup{(\roman*)},leftmargin=*]
\item $\varphi_{a,b}$, if $a\le0$, $b>0$ is nonintegral,
      and $a+b>0$;
\item $\Phi_{a,b,c}$, if $a,b>0$ are nonintegral, $c\le0$,
      and $a+b+c>0$.
\end{enumerate}
\end{lemma}

\begin{proof}
The identities
\begin{align}
 d\bigl(z^a(1-z)^b\bigr)
   &=a\varphi_{a,b}-(a+b)\varphi_{a+1,b},
       \label{eq:curve-contiguity}\\
 d\bigl(X^{a-1}Y^{b-1}W^c(XdY-YdX)\bigr)
   &=(a+b+c)\Phi_{a,b,c+1}-c\Phi_{a,b,c}
       \label{eq:surface-contiguity}
\end{align}
hold on the same cover as the original forms. If the
nonpositive parameter is nonintegral, iterate until it first
becomes positive. The resulting cohomology multiplier is
$(\sigma)_u/(v)_u$, where $v$ is the initial parameter,
$\sigma$ is the initial sum, and $(x)_u=x(x+1)\cdots(x+u-1)$.
It is nonzero. Lemma~\ref{lem:positive} applies to the final
form.

If $v=-m$ is an integer, stop at zero. The multiplier
is $(-1)^m(\sigma)_m/m!$. For the curve form at $a=0$, a
small circle around $z=0$ has nonzero integral, since the
residue is one. For the surface form at $c=0$, the residue
along $W=0$ is, up to sign,
$X^{a-1}(1-X)^{b-1}dX$. Its compact period is nonzero by
Lemma~\ref{lem:positive}. A small normal circle tube over
that cycle pairs with $\Phi_{a,b,0}$ by $2\pi i$ times this
period. The normal character is trivial, so the tube closes.
The nonzero multipliers prove both assertions.
\end{proof}

Integer changes of a parameter multiply forms by
single-valued functions, so all these identities take place
on one cover. If a geometric cover is larger than the
minimal cover of a form, transfer along the finite \'etale
map lifts a cycle and multiplies its integral by the degree.
We will use this to pass between the explicit Euler covers
and the normal covers of divisors.

\section{The cyclic curve cover and the uniform factor}
\label{sec:uniform}

On a finite cover, the residue over $C$ separates into an
Euler form in the fiber and a differential on a compact
curve. Ramification controls the poles of the curve
differential, and Riemann--Roch supplies the numerators
that give its nonzero periods.

\subsection{A divisor over the tangent curve}

Blow up the origin. Near $C$, use a radial coordinate $u_i$,
the normal coordinate $v_i=g_i$, and a curve coordinate
$w_i$. With $h_i^0=h_i(0,w_i)$, the pullbacks are
\begin{equation}\label{eq:first-blowup}
 \begin{split}
 f\circ\pi_0
   &=u_i^A\bigl(v_i^r+u_i^kh_i(v_i,w_i)
                      +u_i^{k+1}Q_i(u_i,v_i,w_i)\bigr),\\
 \pi_0^*\omegao
   &=u_i^2J_i(v_i,w_i)\,du_i\wedge dv_i\wedge dw_i.
 \end{split}
\end{equation}
Set $\eta_i=J_i(0,w_i)dw_i$. Along $C$, a projective
change of scale $\lambda$ gives
\begin{equation}\label{eq:curve-transitions}
 u_b=\lambda u_a,\quad v_b=\lambda^{-e}v_a,\quad
 h_b^0=\lambda^{-D}h_a^0,\quad
 \eta_b=\lambda^{e-3}\eta_a.
\end{equation}
The last relation is the adjunction formula
$K_C\simeq\OO_C(e-3)$ in these local frames.

Let $\gamma=\gcd(r,k)$, $r'=r/\gamma$, $k'=k/\gamma$.
Over $C^\circ=C\setminus P$, where $P$ is the zero set of
$h_D|_C$, take a toroidal modification containing the
primitive ray $(r',k')$ in the $(u,v)$ plane. The two
Cartier divisors $u=0$, $v=0$ make these modifications
compatible on overlaps. A regular subdivision and a
resolution away from the open part under consideration
give a divisor $E_R$ with
\begin{equation}\label{eq:R-numerics}
 N_R=\frac{rD}{\gamma},\qquad
 \nu_R=\frac{3r+k}{\gamma},\qquad
 \frac{\nu_R}{N_R}=\alpha_R.
\end{equation}
The terms containing $h_i-h_i^0$ or $Q_i$ have strictly
larger weight than $N_R$.

\subsection{Separating the curve and the Euler fiber}

Let $p:\widetilde C\to C$ be the compact normalized cyclic
cover defined in the line bundle $\OO_C(1)$ by
\begin{equation}\label{eq:root-curve}
 H_i^D=h_i^0,\qquad H_b=\lambda^{-1}H_a.
\end{equation}
The cover has degree $D$ and may be disconnected. Above
$C^\circ$ it is \'etale. On this cover put
\[
 U=u_iH_i,\qquad V=v_i/H_i^e.
\]
The leading function is $U^A(V^r+U^k)$, and the leading
volume form is
\begin{equation}\label{eq:trivialized-volume}
 H_i^{e-3}U^2dU\wedge dV\wedge p^*\eta_i.
\end{equation}
Choose integers $L,M$ with $r'M-k'L=1$. A unimodular
torus chart at $E_R$ is
\begin{equation}\label{eq:R-chart}
 U=t^{r'}Z^L,\qquad V=t^{k'}Z^M.
\end{equation}
Consequently,
\begin{equation}\label{eq:R-leading}
 \begin{split}
 f\circ\pi&=t^{N_R}Z^{LD}(1+Z^\gamma)
                   +\text{higher powers of }t,\\
 \pi^*\omegao&=t^{\nu_R-1}Z^{3L+M-1}dt\wedge dZ
                   \wedge H_i^{e-3}p^*\eta_i
                   +\text{higher powers of }t.
 \end{split}
\end{equation}

The leading coordinates in \eqref{eq:R-chart} give
the gluing of the normal cover. Extending
$H_i$ locally off $C$ gives on overlaps
$U_b/U_a=1+O(v_a)$ and $V_b/V_a=1+O(v_a)$. Since
$t=U^MV^{-L}$ and $Z=U^{-k'}V^{r'}$, it follows that
\begin{equation}\label{eq:normal-gluing}
 t_b/t_a=1+O(t_a^{k'}),\qquad Z_b/Z_a=1+O(t_a^{k'}).
\end{equation}
Thus $Z$ and the first normal coordinate agree along the
divisor after pullback to $\widetilde C^\circ=p^{-1}(C^\circ)$.
For
\begin{equation}\label{eq:Sigma}
 \Sigma=\{(Z,\tau):Z\ne0,\ Z^\gamma\ne-1,\quad
         \tau^{N_R}Z^{LD}(1+Z^\gamma)=1\},
\end{equation}
we obtain a finite \'etale map of degree $D$
\begin{equation}\label{eq:product-map}
 \mathcal M:\widetilde C^\circ\times\Sigma
                   \longrightarrow\widehat E_R^\circ.
\end{equation}
It is the base-change projection induced by the finite \'etale
cover $\widetilde C^\circ\to C^\circ$.

Let $p_j$ be homogeneous of degree $j$ with nonzero
restriction $a_j=p_j|_C$. For a nonnegative integer $\ell$,
use the original holomorphic numerator $g_e^\ell p_j$.
Its order along $E_R$ is
\begin{equation}\label{eq:numerator-order}
 m_{\ell,j}=\frac{\ell D+jr}{\gamma},\qquad
 \frac{m_{\ell,j}}{N_R}=\frac\ell r+\frac jD.
\end{equation}
Since $g_e=u_i^ev_i=U^eV$, the curve factor of its residue
is independent of $\ell$. Direct substitution in
\eqref{eq:R-leading} gives
\begin{equation}\label{eq:residue-product}
 \mathcal M^*\rho_{g_e^\ell p_j}
     =\phi_{\ell,j}\wedge\theta_j,\qquad
 \theta_j=H_i^{e-3-j}p^*(a_{j,i}\eta_i).
\end{equation}
The transitions \eqref{eq:curve-transitions} and
\eqref{eq:root-curve} show that $\theta_j$ is global.
The fiber factor is the single-valued form on $\Sigma$
\begin{equation}\label{eq:fiber-single-valued}
 \phi_{\ell,j}
   =\tau^{\nu_R+m_{\ell,j}}
      Z^{(3+j+e\ell)L+(\ell+1)M-1}dZ.
\end{equation}
Let $\alpha_{\ell,j}=\alpha_R+\ell/r+j/D$.
Using the defining equation of $\Sigma$ and
$rM-kL=\gamma$, its branch expression reduces to
\begin{equation}\label{eq:fiber-form}
 \phi_{\ell,j}
   =Z^{\gamma(\ell+1)/r-1}(1+Z^\gamma)^{-\alpha_{\ell,j}}dZ
\end{equation}
up to a fixed nonzero constant.
Set $Y=Z^{-\gamma}$ and $z=Y/(1+Y)$. The map
$\Sigma\to M_1$, $z=1/(1+Z^\gamma)$, is finite \'etale.
Under this map,
\eqref{eq:fiber-form} becomes a nonzero multiple of
\begin{equation}\label{eq:fiber-Euler}
 \varphi_{A_j,B_\ell},\qquad
 A_j=\frac{3+j-e}{D},\quad B_\ell=\frac{\ell+1}{r},\quad
 A_j+B_\ell=\alpha_{\ell,j}>0.
\end{equation}
The single-valued expression \eqref{eq:fiber-single-valued}
shows that the Euler character becomes trivial on
$\Sigma$. Thus each connected component maps to the
minimal character cover of $\varphi_{A_j,B_\ell}$.
Transfer to one such component lifts any nonzero period
to $\Sigma$.

\subsection{Ramification supplies the numerators}

The fiber period is nonzero for every pair in
\eqref{eq:triangular-roots}: when $A_j\le0$, use
Lemma~\ref{lem:euler}(i); when $A_j>0$, both $A_j$ and
$B_\ell$ lie in $(0,1)$, so use Lemma~\ref{lem:positive}.
We now construct a nonzero period of $\theta_j$ for every
$0\le j<D$.

For $j\le e-3$, take any nonzero homogeneous $p_j$.
Its restriction to $C$ is nonzero because $j<e$, and
$\theta_j$ is a nonzero holomorphic differential on
$\widetilde C$. It therefore has a nonzero period, which
can be represented away from the finitely many deleted
points.

For $j\ge e-2$, put $v=j+3-e>0$ and write
\[
 \divisor_C(h_D|_C)=\sum_{x\in P}q_xx,\qquad
 \sum_xq_x=eD.
\]
Set $d_x=\gcd(D,q_x)$. On a normalized branch over $x$,
a local parameter $z$ satisfies
$w=z^{D/d_x}\cdot\mathrm{unit}$ and
$H=z^{q_x/d_x}\cdot\mathrm{unit}$. If $a_j$ vanishes to
order $n$ at $x$, then
\begin{equation}\label{eq:ramification-order}
 \ord_z\theta_j=\frac{D(n+1)-vq_x-d_x}{d_x}.
\end{equation}
The ramification contribution $D/d_x-1$ from $p^*\eta_i$
is included here. Because the remainder of $vq_x$ modulo
$D$ is a multiple of $d_x$, this order is nonnegative
exactly when
\begin{equation}\label{eq:floor-condition}
 n\ge\left\lfloor\frac{vq_x}{D}\right\rfloor.
\end{equation}
Define
\begin{equation}\label{eq:RR-line-bundle}
 \Gamma_v=\sum_x\left\lfloor\frac{vq_x}{D}\right\rfloor x,
 \qquad L_j=\OO_C(j)(-\Gamma_v).
\end{equation}
With $g_C=(e-1)(e-2)/2$,
\begin{equation}\label{eq:RR-degree}
 \deg L_j\ge ej-ev=e(e-3)=2g_C-2.
\end{equation}

If $e\ge4$, Riemann--Roch gives
$h^0(C,L_j)\ge g_C-1>0$. A nonzero section $a_j$ produces
a nonzero holomorphic $\theta_j$ on the compact normalized
cover, by \eqref{eq:ramification-order}. It lifts to a
homogeneous polynomial $p_j$: the restriction sequence
\begin{equation}\label{eq:homogeneous-lift}
 0\longrightarrow\OO_{\PP^2}(j-e)
 \longrightarrow\OO_{\PP^2}(j)
 \longrightarrow\OO_C(j)\longrightarrow0
\end{equation}
is surjective on sections, since
$H^1(\PP^2,\OO_{\PP^2}(n))=0$ for every integer $n$.

For $e=3$, the curve is elliptic and $v=j$. A
positive-degree $L_j$ has a section by Riemann--Roch;
the same construction applies whenever $H^0(C,L_j)\ne0$.
It remains to treat $\deg L_j=0$ and
$H^0(C,L_j)=0$. Equality in \eqref{eq:RR-degree} then
forces every $jq_x/D$ to be an integer, say $n_x$.
Since $j>0$, some $n_{x_0}\ge1$. The degree-one bundle
$L_j(x_0)$ has a one-dimensional section space. A
nonzero section cannot belong to $H^0(C,L_j)=0$.
Thus, as a section of $\OO_C(j)$, it vanishes to order
exactly $n_{x_0}-1$ at $x_0$
and at least $n_x$ at the other points. Here
$\Gamma_j-x_0$ is effective, so the section remains
holomorphic on the original curve and lifts by
\eqref{eq:homogeneous-lift}. Formula
\eqref{eq:ramification-order} gives simple poles over
$x_0$ and holomorphy elsewhere. At each such pole its
expansion is $(c_z/z+O(1))dz$ with $c_z\ne0$, so a small
circle has the nonzero period $2\pi i c_z$.

In all cases we have an original polynomial $p_j$ and
a compact cycle $\gamma_j\subset\widetilde C^\circ$ with
$\int_{\gamma_j}\theta_j\ne0$.

\subsection{From product periods to roots}

For a pair $(\ell,j)$ in \eqref{eq:triangular-roots}, choose
the numerator above and a nonzero fiber cycle $\delta_{\ell,j}$
from \eqref{eq:fiber-Euler}. Transfer to $\Sigma$ if
necessary. By \eqref{eq:residue-product}, the pushforward
\[
 \mathcal M_*(\gamma_j\times\delta_{\ell,j})
\]
has nonzero pairing with the specified residue, up to the
fixed orientation sign. Equation \eqref{eq:numerator-order}
gives $m_{\ell,j}<N_R$, so
Proposition~\ref{prop:criterion} yields
$\bfzero(-\alpha_{\ell,j})=0$.
Taking $\ell=0$ gives $D$ distinct roots and hence the
divisibility \eqref{eq:uniform-factor}. This proves
Theorem~\ref{thm:uniform}.

For $r=2$, the triangular set has only $\ell=0$. If
$r\mid k$, then $D/r$ is an integer and the reindexing
$j'=j+\ell D/r$ puts every triangular root in the uniform
block: the window condition gives $0\le j'<D$.
For $r\ge3$ and $r\nmid k$, the allowed pair $(\ell,j)=(1,0)$
gives a root outside that block, since $D/r\notin\ZZ$.

\section{Contact roots}\label{sec:contact}

Assume \eqref{eq:two-layer} and fix $p\in P$ with contact
order $q\ge2$. We construct the contact period explicitly.

\subsection{The local model and its normal cover}

In \eqref{eq:first-blowup}, smoothness of $H$ gives
$h(v,w)=a(v,w)(v+w^q)$ after a change of $w$, where $a$
is a unit. Replacing the radial coordinate by $u a^{1/k}$
puts the function and the original volume form, up to
analytic units, in the form
\begin{equation}\label{eq:contact-model}
 F=u^A(v^r+u^kv+u^kw^q),\qquad
 \omega=u^2du\wedge dv\wedge dw.
\end{equation}
Let
\begin{equation}\label{eq:T-data}
 T'=\bigl(q(r-1),qk,k\bigr),\quad
 t_q=\gcd(k,q(r-1)),\quad T=T'/t_q.
\end{equation}
The compact Newton face is nondegenerate. A regular toric
subdivision containing $T$ gives a divisor $E_T$
with
\begin{equation}\label{eq:T-numerics}
 N_T=qB/t_q,\qquad \nu_T=(q\beta+k)/t_q,\qquad
 \nu_T/N_T=\alpha_q.
\end{equation}
All three weights of $T$ are positive. Thus the analytic
units in the model restrict to nonzero constants on its
divisor torus. Denote the function and volume constants by
$c_f,c_\omega\ne0$, respectively.

To compute that residue, use the finite monomial map
\begin{equation}\label{eq:finite-contact-chart}
 u=t^{q(r-1)}\xi^{-1}\eta^{-(r-1)},\qquad
 v=t^{qk}\eta^{-k},\qquad w=t^k.
\end{equation}
Its exponent matrix has determinant of absolute value $k^2$,
and direct substitution gives
\begin{equation}\label{eq:contact-leading}
 \begin{split}
 F&=t^{qB}\xi^{-D}\eta^{-B}(1+\xi^k+\eta^k),\\
 \omega&=\pm k^2t^{q\beta+k-1}\xi^{-4}\eta^{-\beta-1}
                   dt\wedge d\xi\wedge d\eta.
 \end{split}
\end{equation}
The map factors through a chart of the primitive ray $T$.
On the divisor torus it is finite \'etale, with normal
ramification index $t_q$. A normal coordinate $t_0$ at
$E_T$ satisfies $t_0=t^{t_q}M(\xi,\eta)$ for an invertible
monomial $M$. Define the auxiliary normal cover by
\begin{equation}\label{eq:contact-auxiliary-cover}
 \widehat M_T:\quad
 \vartheta^{qB}\xi^{-D}\eta^{-B}(1+\xi^k+\eta^k)=1,
 \qquad \xi,\eta,\vartheta\ne0.
\end{equation}
Choose $c_1$ with $c_1^{N_T}=c_f^{-1}$. There is a finite
\'etale map $\kappa:\widehat M_T\to\widehat E_T^\circ$
given by $\tau_0=c_1M\vartheta^{t_q}$.
The specified single-valued form on the auxiliary cover is
\begin{equation}\label{eq:contact-single-valued}
 \rho'=\pm k^2\vartheta^{q\beta+k}\xi^{-4}\eta^{-\beta-1}
                         d\xi\wedge d\eta.
\end{equation}
The auxiliary orders are $t_qN_T$ and $t_q\nu_T$.
Using
$dt_0/t_0=t_q\,dt/t+d\log M$, the residue computed in
\eqref{eq:contact-single-valued} satisfies
\begin{equation}\label{eq:contact-residue-comparison}
 \rho'=t_q c_\omega^{-1}c_1^{-\nu_T}\kappa^*\rho_{E_T},
\end{equation}
up to the common orientation sign.

Set $x=\xi^k$, $y=\eta^k$. The factor $k^2$ cancels the
Jacobian of this change, and the residue becomes
\[
 x^{a-1}y^{b-1}(1+x+y)^{-\alpha_q}dx\wedge dy,
 \qquad a=\frac{D\alpha_q-3}{k},\quad
 b=\frac{B\alpha_q-\beta}{k}.
\]
The projective change
$X=x/(1+x+y)$, $Y=y/(1+x+y)$, $W=1/(1+x+y)$ gives
$\Phi_{a,b,c}$ with
\begin{equation}\label{eq:contact-Euler}
 \begin{gathered}
 a=\frac{D-3+D/q}{B}>0,\qquad b=\frac1q,\\
 c=-\frac{r[q(e-3)+D]}{qB}<0,\qquad a+b+c=\alpha_q>0.
 \end{gathered}
\end{equation}
These substitutions define a finite \'etale map
$\widehat M_T\to M_2$. Equation
\eqref{eq:contact-single-valued} shows that it lifts to
the minimal Euler character cover. Transfer followed by
$\kappa_*$ therefore preserves nonvanishing of periods,
by \eqref{eq:contact-residue-comparison}.

\subsection{Proof of the contact root theorem}

If $r\ge3$, then $0<a<1$. If $r=2$, then $0<a<3/2$;
its only possible integral value is $a=1$.
This value occurs exactly when
\begin{equation}\label{eq:integer-resonance}
 r=2,\qquad D=q(k+3),\qquad
 (a,b,c)=(1,1/q,-1),\qquad \alpha_q=1/q.
\end{equation}
The identity $2e=(q-1)k+3q$ forces $k,q$ to be even.
For $j=k/2$,
\[
 e-3-j=\frac{(q-2)(k+3)}2\ge0,\qquad
 \alpha_R+\frac jD=\frac{k+3}{D}=\frac1q.
\]
Theorem~\ref{thm:uniform} gives $\bfzero(-\alpha_q)=0$.
The associated torus form has the explicit primitive
\begin{equation}\label{eq:exact-contact-form}
 \Phi_{1,b,c}=d\left(\frac{Y^bW^{c-1}}b\,dW\right),
 \qquad X=1-Y-W.
\end{equation}

For all other parameters, $a,b>0$ are nonintegral and
$c<0$. Lemma~\ref{lem:euler}(ii) supplies a nonzero
compact Euler period. Transfer followed by $\kappa_*$
gives a nonzero period of $\rho_{E_T}$, and
Proposition~\ref{prop:criterion} with numerator $1$
gives $\bfzero(-\alpha_q)=0$.
This completes the proof of Theorem~\ref{thm:contact}.

\section{Topological poles and strong monodromy}\label{sec:poles}

For the two-layer family \eqref{eq:two-layer}, the
topological zeta function is determined by $e,r,k$ and
the contact orders. We compute it and identify its poles
among the roots constructed above.

These hypotheses make $f$ reduced with an isolated
critical point at the origin. Indeed, if $f=df=0$ at a
nonzero point, the homogeneous Euler identity gives
\[
 Ag_e^r+Dh_D=0,\qquad g_e^r+h_D=0.
\]
Thus $g_e=h_D=0$, where $df=dh_D\ne0$ by smoothness of
$H$ at $P$, a contradiction. A positive-dimensional
critical branch through the origin would have $f$
constant and hence zero, so the critical point is isolated.
A repeated factor would produce a positive-dimensional
critical set, proving reducedness.

Let $(E_i)_{i\in J}$ be the components of the function's
total transform on a resolution $\pi$, including the strict
transforms. Their numerical data are
\[
 N_i=\ord_{E_i}(f\circ\pi),\qquad
 \nu_i-1=\ord_{E_i}(\pi^*\omegao).
\]
For $\varnothing\ne I\subseteq J$, set
\[
 E_I^\circ=\left(\bigcap_{i\in I}E_i\right)
                       \setminus\bigcup_{j\notin I}E_j.
\]
Our convention is
\begin{equation}\label{eq:zeta-definition}
 Z_{\mathrm{top},0}(f;s)=
 \sum_{\varnothing\ne I\subseteq J}
 \chi(E_I^\circ\cap\pi^{-1}(0))
                    \prod_{i\in I}\frac1{N_is+\nu_i}.
\end{equation}
Here $\omegao=dx\wedge dy\wedge dz$.

\subsection{The complete zeta function}

Write $\chi_C=e(3-e)$, $b=|P|$, and set
\begin{equation}\label{eq:zeta-denominators}
 \begin{gathered}
 a=As+3,\qquad d=s+1,\qquad \mathcal R=rDs+3r+k,\\
 K_q=qBs+q\beta+k,\qquad
 \gamma=\gcd(r,k),\quad h=\gcd(k,r-1),\\
 j_q=\gcd(k,r,q),\qquad t_q=\gcd(k,q(r-1)),\\
 G_q=\mathcal R(t_q^2-qh^2)
               -k\bigl(j_q^2+q(r-1)-r\bigr).
 \end{gathered}
\end{equation}

\begin{proposition}[Weighted zeta formula]\label{prop:zeta}
For \eqref{eq:two-layer},
\begin{equation}\label{eq:global-zeta}
 \begin{split}
 Z_{\mathrm{top},0}(f;s)
   ={}&\frac{3-\chi_C}{a}+\frac{k\chi_C}{a\mathcal R}\\
     &+\frac{r\chi_C+(\chi_C-b)(r-\gamma^2)s}
                  {d\mathcal R}
       +\frac{s}{d\mathcal R}
                   \sum_{p\in P}\frac{G_{q_p}}{K_{q_p}}.
 \end{split}
\end{equation}
Consequently, every pole is simple and belongs to
\begin{equation}\label{eq:pole-list}
 \left\{-1,-\frac3A,-\alpha_R\right\}
        \cup\{-\alpha_q:q=q_p\ge2\}.
\end{equation}
\end{proposition}

\begin{proof}
After the first blow-up, the three types of strata have
local models
\[
 u^A\quad\text{on }\PP^2\setminus C,\qquad
 G_k=u^A(v^r+u^k)\quad\text{on }C\setminus P,
\]
and $F_{q,k}=u^A(v^r+u^kv+u^kw^q)$ at a point of contact
order $q$. At a transverse intersection one obtains the
same model with $q=1$ by taking $w=h-v$.
All use the volume weight $u^2$.
The local Newton calculation in Appendix~\ref{app:newton}
gives
\begin{equation}\label{eq:local-zeta-summary}
 \begin{gathered}
 Z_{G_k}=\frac{k/a+r-\gamma^2s/d}{\mathcal R},\qquad
 Z_T=\frac{k/a+r/d}{\mathcal R},\\
 Z_{F_{q,k}}=Z_T+\frac{sG_q}{d\mathcal R K_q}.
 \end{gathered}
\end{equation}
In particular $G_1=0$.
Euler integration over the strata of the first exceptional
plane gives
\begin{equation}\label{eq:zeta-assembly}
 Z_{\mathrm{top},0}(f;s)
   =\frac{3-\chi_C}{a}+(\chi_C-b)Z_{G_k}
                         +\sum_{p\in P}Z_{F_{q_p,k}}.
\end{equation}
This is additivity of the finite sum
\eqref{eq:zeta-definition}, followed by Euler integration
along the map to the first exceptional divisor. Substituting
\eqref{eq:local-zeta-summary} proves \eqref{eq:global-zeta}.

The denominator gives \eqref{eq:pole-list}. For the
relative positions of its values, observe that
\begin{equation}\label{eq:slope-order}
 \alpha_R-\frac3A=\frac{k(e-3)}{erD},\qquad
 \alpha_q-\alpha_R=
          \frac{k[q(e-3)+D]}{qDB}>0.
\end{equation}
The numbers $\alpha_q=\beta/B+k/(qB)$ are distinct for
distinct $q\ge2$, and $\alpha_R<\alpha_q<1$.
For $e\ge4$ all entries of \eqref{eq:pole-list} are
distinct. For $e=3$, we have
$a=3(rs+1)$ and $\mathcal R=D(rs+1)$; the only term
that could have a double pole is
$k\chi_C/(a\mathcal R)$, which vanishes since $\chi_C=0$.
Thus every pole is simple.
\end{proof}

\subsection{The first-divisor root}

The first exceptional divisor gives the root $-3/A$
throughout the analytic family \eqref{eq:analytic-family}.

\begin{proposition}\label{prop:first-root}
Every germ \eqref{eq:analytic-family} satisfies
$\bfzero(-3/A)=0$.
\end{proposition}

\begin{proof}
The first exceptional divisor has open part $\PP^2\setminus C$
and data $(A,3)$. One component of its normal cover is
$\tau^eg_i=1$. Setting $T_i=\tau^{-1}$ identifies it with
$U=S\setminus B_0$, where
\[
 S=\{T^e=g_e(X,Y,Z)\}\subset\PP^3,\qquad B_0=\{T=0\}\simeq C.
\]
The projection $p_0:S\to\PP^2$ is finite, and $S$ is
smooth. With $\Omega_i$ the projective chart area form,
its specified residue is $\sigma=p_0^*\Omega_i/T_i^3$.
The numerator and denominator both change by
$\lambda^{-3}$ on overlaps, so $\sigma$ is global.
Near $B_0$, write
$\Omega_i=J(v,w)dv\wedge dw$ and $v=T_i^e$. Then
\begin{equation}\label{eq:cyclic-surface-form}
 \sigma=eT_i^{e-4}J(T_i^e,w)dT_i\wedge dw.
\end{equation}
For $e\ge4$, this is a nonzero holomorphic two-form on
the compact surface $S$, with
$\int_S\sigma\wedge\overline\sigma>0$. If its class
vanished on $U$, the localization sequence
\[
 H^0(B_0,\CC)\longrightarrow H^2(S,\CC)
                         \longrightarrow H^2(U,\CC)
\]
would express $[\sigma]$ as a linear combination of the
Poincar\'e duals of the components of $B_0$. Pairing with
$\overline\sigma$ gives zero, since its restriction to
each curve is zero, a contradiction. Thus $\sigma$ has a
nonzero compact period in $U$.
For $e=3$, \eqref{eq:cyclic-surface-form} is logarithmic
with residue $3\eta$, where $\eta$ is a nonzero holomorphic
differential on the elliptic curve $C$. A normal circle tube
over a nonzero period of $\eta$ gives the required period.
Proposition~\ref{prop:criterion} with numerator $1$ finishes
the proof.
\end{proof}

\subsection{From poles to Bernstein--Sato roots}

By Proposition~\ref{prop:zeta}, each actual pole is listed
in \eqref{eq:pole-list}. Proposition~\ref{prop:first-root}
handles $-3/A$, Theorem~\ref{thm:uniform} with $j=0$
handles $-\alpha_R$, and Theorem~\ref{thm:contact} handles
the contact values. Finally, $\bfzero(-1)=0$: substituting
$s=-1$ in a Bernstein identity makes its right side
$P(-1)1$ holomorphic, whereas $b(-1)/f$ cannot be
holomorphic unless $b(-1)=0$. This proves
Corollary~\ref{cor:strong}.

\section{Roots beyond the pole set}\label{sec:examples}

In the family below, a contact root has a nonzero period
even though its topological candidate cancels. We then
exhibit an actual pole detected by the curve period.

\subsection{An infinite family with a cancelled contact pole}

Assume \eqref{eq:cancelled-family}. Fix a smooth degree-$e$
curve $C$ and a point $x_0\in C$. There is a section of
$L=\OO_C(D)$ with a zero of order exactly $q$ at $x_0$
and only simple zeros elsewhere. Indeed,
\[
 \deg L-q-2=eD-q-2>2g_C-2.
\]
Riemann--Roch makes one additional vanishing condition
at $x_0$ a proper hyperplane in $H^0(C,L(-qx_0))$.
For each $x\ne x_0$, a double zero at $x$ imposes two
independent conditions. The incidence space of these
sections as $x$ varies on a curve has codimension at least
one. A general section has the stated divisor.
Lift it to $h_D$ using \eqref{eq:homogeneous-lift}.
Adding $g_e a_{D-e}$ adjusts the normal derivative at
$x_0$ freely, so $H$ can be chosen smooth there.
At the other zeros, simplicity along $C$ already implies
smoothness of $H$. This proves geometric existence.

For these parameters,
\[
 B=k(4q+1),\qquad \beta=4k,\qquad
 \mathcal R=4rq s+4r-1,\qquad
 K_q=k(4q+1)(qs+1).
\]
The lattice indices are $h=t_q=k$ and $j_q=1$. Hence
\begin{equation}\label{eq:cancelled-G}
 G_q=-k^2(q-1)(\mathcal R+1),\qquad
 \frac{G_q}{K_q}=-\frac{4rk(q-1)}{4q+1}.
\end{equation}
The contact term in \eqref{eq:global-zeta} has lost its
$K_q$ denominator. At $s=-1/q$ the other denominators
satisfy
\[
 a=\frac{r-1-q}{q}\ne0,\qquad
 \mathcal R=-1,\qquad d=1-1/q\ne0;
\]
here $q\ge r$. All remaining contacts are transverse and
have $G_1=0$. Thus the complete zeta function is
holomorphic at $-1/q$.

The primitive contact ray is
$T=(q,q,1)$, with
\[
 (N_T,\nu_T)=(q(4q+1),4q+1).
\]
Its prescribed Euler parameters are
\[
 (a,b,c)=\left(\frac1{r-1},\frac1q,-\frac1{r-1}\right).
\]
Equation \eqref{eq:surface-contiguity} gives on the same cover
\begin{equation}\label{eq:cancelled-nonzero}
 [\Phi_{1/(r-1),1/q,-1/(r-1)}]
   =-\frac{r-1}{q}
      [\Phi_{1/(r-1),1/q,1-1/(r-1)}].
\end{equation}
All three parameters on the right are positive and
nonintegral, so its compact period is nonzero.
The comparison \eqref{eq:contact-residue-comparison}
gives a nonzero period of the contact residue, and
Proposition~\ref{prop:criterion} gives $\bfzero(-1/q)=0$.

Finally,
\[
 \frac1q-\alpha_R=\frac1{rD}.
\]
Membership in \eqref{eq:triangular-roots} would require
$\ell D+rj=1$ for nonnegative integers $\ell,j$, which
is impossible. This proves Theorem~\ref{thm:cancelled}.

\begin{example}[At least thirty-seven distinct roots]
\label{ex:thirty-seven}
Take $(e,r,k,q)=(6,3,2,5)$ and any pair constructed above.
The triangular roots are
\[
 \left\{-\frac{11+3j}{60}:0\le j\le19\right\}
 \quad\cup\quad
 \left\{-\frac{31+3j}{60}:0\le j\le13\right\}.
\]
The two numerator progressions differ modulo three, so
there are $34$ values. The roots $-1/6$, $-1/5$, and
$-1$ lie outside this set. Thus $\deg\bfzero\ge37$,
although $-1/5$ is absent from the topological pole set.
\end{example}

\subsection{An actual pole with an exact contact torus form}

The following integral-coefficient example realizes the
resonance \eqref{eq:integer-resonance}. Put
\begin{equation}\label{eq:integer-example}
 \begin{gathered}
 g=x^4+y^4+z^4,\quad p_8=x^8+y^8+z^8,\quad
 p_6=x^6+y^6+z^6,\\
 f=g^2+z^2p_8+gp_6.
 \end{gathered}
\end{equation}
Here $(e,r,k,D)=(4,2,2,10)$. The four points of
$C\cap\{z=0\}$ have contact order two, since $p_8\ne0$
there. At these points $p_6=y^4(y^2-x^2)\ne0$, so
$H=V(z^2p_8+gp_6)$ is smooth. The remaining intersections
are the $32$ transverse points of $C\cap V(p_8)$.
To check transversality, all coordinates at these points
are nonzero, and the \'etale map
$[x:y:z]\mapsto[x^4:y^4:z^4]$ reduces the intersection
to the transverse line and conic
$U+V+W=0$, $U^2+V^2+W^2=0$.
Thus $n_2=4$, $n_1=32$, $b=36$, and $\chi_C=-4$.

With $a=8s+3$, $\mathcal R=20s+8$, $d=s+1$, and
$K_2=24s+12$, formula \eqref{eq:global-zeta} becomes
\begin{equation}\label{eq:example-zeta}
 Z_{\mathrm{top},0}(f;s)
   =\frac7a-\frac8{a\mathcal R}
       +\frac{80s-8}{d\mathcal R}
       +\frac{4s(40s+8)}{d\mathcal R K_2}.
\end{equation}
Its pole at $-1/2$ has residue $-1$. The corresponding
contact torus form has parameters $(1,1/2,-1)$ and is
exact by \eqref{eq:exact-contact-form}.
On $E_R$, the numerator $p=x$ has order $r/\gamma=1$,
with $\gamma=\gcd(r,k)=2$ and $(N_R,\nu_R)=(10,4)$.
Its fiber form has parameters
$(0,1/2)$, and its curve form is the nonzero holomorphic
differential $p^*(x_i\eta_i)$. These two periods detect
the exact root $-(4+1)/10=-1/2$.
Theorem~\ref{thm:uniform} gives in total the ten roots
\[
 -\frac25,-\frac12,-\frac35,-\frac7{10},-\frac45,
 -\frac9{10},-1,-\frac{11}{10},-\frac65,-\frac{13}{10};
\]
the first-divisor root $-3/8$ is an additional one.

\subsection{A degree bound from distinct roots}

The uniform factor, the first-divisor root and the
distinct contact roots give a lower bound for the degree.

\begin{corollary}\label{cor:degree}
For every germ \eqref{eq:analytic-family},
\begin{equation}\label{eq:degree-bound}
 \deg\bfzero\ge D+\mathbf1_{e\ge4}+\mathbf1_{r\nmid k}.
\end{equation}
Under \eqref{eq:two-layer}, let $Q$ be the set of distinct
contact orders $q\ge2$ and put
\[
 \delta_q=D(\alpha_q-\alpha_R)
          =\frac{k[q(e-3)+D]}{qB}.
\]
Then the right side of \eqref{eq:degree-bound} can be
increased by $\#\{q\in Q:\delta_q\notin\ZZ\}$.
\end{corollary}

\begin{proof}
For $e\ge4$, the first-divisor root lies above the entire
uniform block by \eqref{eq:slope-order}. The block contains
$-1$ exactly when $D(1-\alpha_R)=D-3-k/r$ is an integer,
equivalently $r\mid k$. Finally,
$\alpha_R<\alpha_q<1$ implies that the contact root is
in the block exactly when $\delta_q$ is an integer.
Distinct contact orders have distinct roots.
\end{proof}

For $k=1$, one has $0<\delta_q<1$ for every actual
$q\ge2$. Thus every distinct nontransverse contact order
adds a root outside the uniform factor.

\appendix
\section{The weighted Newton calculation}\label{app:newton}

This appendix proves \eqref{eq:local-zeta-summary}, keeping
the volume factor $u^2$ from the first blow-up. The
notation is that of \eqref{eq:zeta-denominators}.

\subsection{Faces, cones, and lattice indices}

The Newton vertices of
$F_{q,k}=u^A(v^r+u^kv+u^kw^q)$ are
\[
 V_0=(A,r,0),\qquad V_1=(D,1,0),\qquad V_2=(D,0,q).
\]
Besides the coordinate normals, use the integer rays
\[
 \mathsf R=(r,k,0),\quad \mathsf Q=(0,q,1),\quad
 \mathsf T=(q(r-1),qk,k).
\]
For a lattice vector $v$, the numerical functions are
\[
 N(v)=\min_i\langle v,V_i\rangle,\qquad
 \nu(v)=3v_1+v_2+v_3.
\]
Thus their linear forms $sN+\nu$ are $\mathcal R$,
$P_q=qs+q+1$, and $K_q$.
Ray rescaling multiplies the cone multiplicity and its
denominator by the same factor. Thus the displayed rays
compute the cone quotients directly; divisor data use
primitive generators, as in \eqref{eq:T-numerics}.

All compact faces are nondegenerate on their tori.
The edges are binomials with normalized lengths
$h,j_q,1$, respectively. The triangle has normalized
area $t_q$. Its curve maps by an isogeny of degree
$t_q$ to $1+x+y=0$, so its Euler characteristic is
$-t_q$. A regular Newton subdivision therefore gives
an actual resolution over the local model.

The weighted Denef--Loeser formula~\cite{DL92} is obtained
from the toric orbits as follows.
For a simplicial cone $\Delta$ with primitive generators
$v_i$, write
$J_\Delta=\mult(\Delta)/\prod_i(sN(v_i)+\nu(v_i))$.
This expression is additive under subdivisions: on
the saturated span it is the cone integral of
$\exp(-sN-\nu)$. The coordinate denominators for the
standard basis vectors $e_1,e_2,e_3$ are $a,1,1$.
The coordinate boundaries for $v=0$ and $w=0$ have
$(N,\nu)=(0,1)$; they refine the orbit stratification
and contribute denominator factor $1$.
Write $J_\tau$ for the cone sum attached to a compact
face $\tau$. The orbit calculation gives
\begin{equation}\label{eq:weighted-Newton-sum}
 Z_{F_{q,k}}=\sum_{\dim\tau=0}J_\tau
       +\frac{s}{d}\sum_{\dim\tau\ge1}
          (-1)^{\dim\tau}\operatorname{Vol}_{\ZZ}(\tau)J_\tau.
\end{equation}
Here normalized volume is lattice length for an edge
and twice lattice area for a triangle. Indeed, the face
hypersurface $V_\tau$ in its effective torus has
$\chi(V_\tau)=(-1)^{\dim\tau-1}\operatorname{Vol}_{\ZZ}(\tau)$.
The torus has Euler characteristic zero. Its complement
and its intersection with the strict transform therefore
give the combined coefficient
\[
 -\chi(V_\tau)+\frac{\chi(V_\tau)}d
    =(-1)^{\dim\tau}\operatorname{Vol}_{\ZZ}(\tau)\frac{s}{d}.
\]
A vertex contributes one. Orbits with an extra torus
factor have Euler characteristic zero. This proves
\eqref{eq:weighted-Newton-sum} with the original
volume weight.

The vertex cones admit the decompositions
\[
 (e_1,e_3,\mathsf T),(e_1,\mathsf T,\mathsf R);\qquad
 (e_3,\mathsf Q,\mathsf T);\qquad
 (e_2,\mathsf R,\mathsf T),(e_2,\mathsf T,\mathsf Q).
\]
Their determinants have absolute values
$qk,k^2,q^2(r-1),rk,q(r-1)$ in that order.
For the edge cones, the gcds of the cross-product
coordinates are $qh,kj_q,q(r-1)$; the ray $\mathsf T$
has lattice index $t_q$. Multiplying by the face
volumes gives the following table. Only the edge
and triangle rows are to be multiplied by $s/d$.
\begin{center}
\renewcommand{\arraystretch}{1.5}
\begin{tabular}{@{}ll@{}}
\toprule
Face & Contribution before the edge/triangle factor $s/d$\\
\midrule
$V_0$ & $\displaystyle\frac{qk}{aK_q}
                       +\frac{k^2}{a\mathcal R K_q}$\\
$V_1$ & $\displaystyle\frac{q^2(r-1)}{P_qK_q}$\\
$V_2$ & $\displaystyle\frac{rk}{\mathcal R K_q}
                       +\frac{q(r-1)}{P_qK_q}$\\
$V_0V_1$ & $\displaystyle-\frac{qh^2}{K_q}$\\
$V_0V_2$ & $\displaystyle-\frac{kj_q^2}{\mathcal R K_q}$\\
$V_1V_2$ & $\displaystyle-\frac{q(r-1)}{P_qK_q}$\\
$V_0V_1V_2$ & $\displaystyle\frac{t_q^2}{K_q}$\\
\bottomrule
\end{tabular}
\end{center}

\subsection{Local simplification and the noncompact candidate}

The $P_q$ terms combine as
\[
 \frac{q(r-1)}{P_qK_q}\left(q+1-\frac{s}{d}\right)
          =\frac{q(r-1)}{dK_q}.
\]
Consequently,
\begin{equation}\label{eq:local-F}
 \begin{split}
 Z_{F_{q,k}}=\frac1{K_q}\Bigl[{}&
       \frac{qk}{a}+\frac{k^2}{a\mathcal R}
       +\frac{rk}{\mathcal R}+\frac{q(r-1)}d\\
       &+\frac{s}{d}\left(t_q^2-qh^2
                          -\frac{kj_q^2}{\mathcal R}\right)\Bigr].
 \end{split}
\end{equation}
The two-dimensional version of the same calculation
for $G_k=u^A(v^r+u^k)$ has vertex terms
$k/(a\mathcal R)$ and $r/\mathcal R$, and edge term
$-\gamma^2s/(d\mathcal R)$. Hence
\begin{equation}\label{eq:local-G}
 Z_{G_k}=\frac{k/a+r-\gamma^2s/d}{\mathcal R}.
\end{equation}

For $q=1$ one has $j_1=1$, $t_1=h$, and
\eqref{eq:local-F} reduces to
$Z_T=(k/a+r/d)/\mathcal R$. For arbitrary $q$, use
\[
 K_q=q(\mathcal R-a)+k,\qquad
 \mathcal R-ra=k(rs+1)
\]
in \eqref{eq:local-F}. Common denominators give
\begin{equation}\label{eq:local-replacement}
 Z_{F_{q,k}}=Z_T+
 \frac{s\{\mathcal R(t_q^2-qh^2)
            -k(j_q^2+q(r-1)-r)\}}{d\mathcal R K_q}.
\end{equation}
Equations \eqref{eq:local-G} and \eqref{eq:local-replacement}
are exactly \eqref{eq:local-zeta-summary}. For $k=1$,
all displayed gcds equal one and
$G_q=-(q-1)r(Ds+4)$, recovering the contact replacement
formula of~\cite{Jiang26}.

\section*{Acknowledgments}

This mathematical paper was prepared with the assistance of artificial intelligence.  I
first acknowledge the textual data produced by humanity as a whole, which
provides the linguistic and cultural record from which mathematical expression
can be learned.  I then thank the accumulated scholarly achievements of
mathematicians throughout history.  I also thank the researchers whose work has
advanced and refined neural networks through successive generations.  Finally,
I thank the OpenAI team for the training and development of the models used in
the preparation of this paper.

\begin{thebibliography}{99}

\bibitem{AGGL26}
E.~Artal Bartolo, P.~D.~Gonz\'alez P\'erez,
M.~Gonz\'alez Villa, and E.~Le\'on Cardenal,
\emph{Denef--Loeser zeta functions of suspensions and
L\^e--Yomdin singularities},
\href{https://arxiv.org/abs/2604.24523v1}{arXiv:2604.24523v1}
(2026).

\bibitem{BV26}
D.~Bath and W.~Veys,
\emph{The Strong Monodromy Conjecture for a class of
homogeneous polynomials in three variables},
\href{https://arxiv.org/abs/2602.20922v1}{arXiv:2602.20922v1}
(2026).

\bibitem{Bla26}
G.~Blanco,
\emph{Topological roots of the Bernstein--Sato polynomial
of plane curves},
\href{https://arxiv.org/abs/2406.09034v3}{arXiv:2406.09034v3}
(2026).

\bibitem{DL92}
J.~Denef and F.~Loeser,
\emph{Caract\'eristiques d'Euler--Poincar\'e, fonctions z\^eta
locales et modifications analytiques},
J. Amer. Math. Soc. \textbf{5} (1992), no.~4, 705--720.

\bibitem{Est26}
A.~Esterov,
\emph{The monodromy conjecture for high-order poles,
and norms of residue antiderivatives},
\href{https://arxiv.org/abs/2610.22111v1}{arXiv:2610.22111v1}
(2026).

\bibitem{GLM97}
S.~M.~Gusein-Zade, I.~Luengo, and A.~Melle-Hern\'andez,
\emph{Partial resolutions and the zeta-function of a singularity},
Comment. Math. Helv. \textbf{72} (1997), no.~2, 244--256.

\bibitem{Jiang26}
Y.~Jiang,
\emph{Contact Reconstruction and Spectral Balance for Surface
Singularities with a Multiple Smooth Tangent Cone},
manuscript, 19 October 2026.

\end{thebibliography}
\end{document}
